\section{Proofs of the sharp global asymptotic and its converses}
\label{app:sharp-global}

Throughout this appendix the assumptions are the non-acceleration
hypotheses of Theorem~\ref{thm:distribution_order}: $p_0\in\mathcal P_2(\R^d)$,
$u\in C^2$, $u_t$ is globally $L$-Lipschitz with uniformly bounded
$u_t(0)$, $\acc_t$ is globally $L_a$-Lipschitz, and the stated uniform
$L^2(p_t)$ bounds on acceleration and material jerk hold.  The
acceleration condition is not imposed.  All constants below may depend
on this fixed field and on $T$, but not on the grid.  The proof does not
assume a positive density, a Poincar\'e inequality, a fourth moment, or
a smooth weighted-gradient projection.

\subsection{Reference coordinates and transport of solenoidal fields}

For $z\in\R^d$, set
\begin{align*}
 X_t(z)&=F_{0,t}(z),& J_t(z)&=DF_{0,t}(z),\\
 c_t(z)&=J_t(z)^{-1}\acc_t(X_t(z)),&
 C_t^\sigma(z)&=\int_0^t\sigma(s)c_s(z)\,\dd s.
\end{align*}
Then
\begin{equation}
 B_t^\sigma(X_t(z))=J_t(z)C_t^\sigma(z).
 \label{eq:global-reference-defect}
\end{equation}
Standard flow bounds, used here directly, give constants $C<\infty$
such that, uniformly in $t\in[0,T]$ and $z$,
\begin{align}
 |X_t(z)|+|\acc_t(X_t(z))|&\le C(1+|z|),\nonumber\\
 \|J_t(z)\|_{\rm op}+\|J_t(z)^{-1}\|_{\rm op}&\le2e^{LT}.
 \label{eq:global-flow-growth}
\end{align}
Indeed $\sup_t|\acc_t(0)|<\infty$ follows from continuity and compactness
in time, so the assumed Lipschitz bound gives linear growth of
acceleration.  The spatial derivative satisfies
$\partial_tJ_t=Du_t(X_t)J_t$.  In particular, as multiplication
operators on $\mathcal H=L^2(p_0;\R^d)$,
\begin{equation}
 \|J_t-J_r\|_{\mathcal H\to\mathcal H}
 \le Le^{LT}|t-r|.
 \label{eq:jacobian-operator-continuity}
\end{equation}
Moreover $c_t$ is continuous in $\mathcal H$ by dominated convergence.
In fact it is Lipschitz there: differentiating along exact trajectories
gives
\[
 \partial_tc_t
 =J_t^{-1}\{\jerk_t(X_t)-Du_t(X_t)\acc_t(X_t)\},
\]
whose $\mathcal H$ norm is at most $e^{LT}(C_j+LC_a)$.  Thus
$C_t^\sigma$ is Lipschitz in $\mathcal H$ for every bounded measurable
$\sigma$.  The field $B_t^\sigma$ is continuous in $(t,x)$, locally
$C^1$ in $x$ with spatial derivatives bounded uniformly on compact
sets for $0\le t\le T$, and has uniformly linear growth.  These claims
follow by differentiating the reference formula locally: $u\in C^2$
makes $J_t$ locally $C^1$ in $z$, $\acc\in C^1$, and all flow images
and inverse images of a fixed compact set remain in a common compact
set.  Differentiation under the bounded-time integral is then justified
by its local continuous bounds.

\begin{lemma}[Transport of the weighted solenoidal space]
\label{lem:global-solenoidal-transport}
Let $F$ be a $C^1$ bi-Lipschitz diffeomorphism, $q=F_\#p$, and define
\[
 (F_*w)(y)=DF(F^{-1}y)w(F^{-1}y).
\]
This is a bounded invertible map from $L^2(p)$ to $L^2(q)$ and maps
$\mathcal G_p^\perp$ onto $\mathcal G_q^\perp$.
\end{lemma}

\begin{proof}
The norm statements follow from the derivative bounds for $F$ and its
inverse.  For a compactly supported smooth test $\phi$,
\[
 \int\nabla\phi\cdot F_*w\,\dd q
 =\int\nabla(\phi\circ F)\cdot w\,\dd p.
\]
The composition is compactly supported and $C^1$.  Such tests can be
approximated by smooth compactly supported tests with uniformly
convergent gradients, so the right side vanishes if
$w\in\mathcal G_p^\perp$.  Apply the same argument to $F^{-1}$ for
the converse.
\end{proof}

\begin{lemma}[Continuity of the projected global coefficient]
\label{lem:global-coefficient-continuity}
For bounded $\sigma$, the function
$K(t)=\|\Pi_{p_t}B_t^\sigma\|_{L^2(p_t)}$ is Lipschitz on $[0,T]$.
\end{lemma}

\begin{proof}
Put $S=\mathcal G_{p_0}^\perp$, a fixed closed subspace of $\mathcal H$.
Lemma~\ref{lem:global-solenoidal-transport} and
Eq.~\eqref{eq:global-reference-defect} give
\begin{equation}
 K(t)=\inf_{z\in S}\|J_t(C_t^\sigma-z)\|_{\mathcal H}.
 \label{eq:global-fixed-hilbert-distance}
\end{equation}
The norms $q_t(v)=\|J_tv\|_{\mathcal H}$ are uniformly equivalent to
the norm of $\mathcal H$: with $m=e^{-LT}$ and $M=e^{LT}$,
$m\|v\|\le q_t(v)\le M\|v\|$.  The infimum has a minimizer in
$S$ because $q_t$ is an equivalent Hilbert norm.  If $z_r$ is a
minimizer at $r$, then
\begin{align*}
 K(t)-K(r)
 &\le M\|C_t^\sigma-C_r^\sigma\|_{\mathcal H}
   +\|J_t-J_r\|_{\mathcal H\to\mathcal H}
       \|C_r^\sigma-z_r\|_{\mathcal H}\\
 &\le M\|C_t^\sigma-C_r^\sigma\|_{\mathcal H}
   +m^{-1}\|J_t-J_r\|_{\mathcal H\to\mathcal H}K(r).
\end{align*}
The same bound with $t,r$ exchanged controls the other sign.
$C_t^\sigma$ and $J_t$ are Lipschitz in the indicated norms and
$K(t)\le M\sup_r\|C_r^\sigma\|_{\mathcal H}<\infty$.  This proves
the assertion without differentiating $\Pi_{p_t}$.
\end{proof}

\subsection{A tangent expansion uniform over the probability path}

The next lemma supplies the exact norm needed after the classical
trajectory expansion.  Its single-measure version is the standard
Wasserstein tangent norm~\citep[Chapter~8]{ambrosio2005gradient}; the
proof below records the uniformity and domination needed here.

\begin{lemma}[Uniform tangent expansion]
\label{lem:uniform-global-tangent}
Let $p_t=(X_t)_\#p_0$ be the flow family above.  Suppose $b_t(x)$ is
continuous in $(t,x)$, locally $C^1$ in $x$ with derivatives bounded
uniformly on compact sets in time and space, and
$|b_t(x)|\le C(1+|x|)$.  Suppose also that
$k(t)=\|\Pi_{p_t}b_t\|_{L^2(p_t)}$ is continuous.  Then
\begin{equation}
 \sup_{t\in[0,T]}
 \left|\frac{W_2((\mathrm{Id}+\epsilon b_t)_\#p_t,p_t)}{|\epsilon|}
       -k(t)\right|\longrightarrow0
 \quad\text{as }\epsilon\to0.
 \label{eq:uniform-tangent-expansion}
\end{equation}
\end{lemma}

\begin{proof}
The uniform linear-growth bounds imply
$M_b=\sup_t\|b_t\|_{L^2(p_t)}<\infty$ and uniform tail control.
Explicitly, $|X_t(z)|\le C_0(1+|z|)$ implies
\begin{align*}
 &\sup_t\int_{|x|>R}|b_t(x)|^2\,p_t(\dd x)\\
 &\qquad\le C_1\int_{|z|>R/C_0-1}(1+|z|)^2\,p_0(\dd z)
 \longrightarrow0.
\end{align*}
Let $\chi_R$ be a smooth spatial cutoff, equal to one on the ball of
radius $R$ and zero outside the ball of radius $2R$, and put
$b_t^R=\chi_Rb_t$.  For fixed $R$ the fields $b_t^R$ are globally
Lipschitz with a common constant $L_R$ and
\[
 d_R:=\sup_t\|b_t-b_t^R\|_{L^2(p_t)}\longrightarrow0.
\]
Apply Lemma~\ref{lem:projected-frozen-flow} to $b_t^R$ or $-b_t^R$
according to the sign of $\epsilon$.  The difference between its frozen
flow and the map $\mathrm{Id}+\epsilon b_t^R$ has $L^2(p_t)$ norm
at most
$L_R\epsilon^2e^{L_R|\epsilon|}\|b_t^R\|_{L^2(p_t)}/2$.
Coupling the two Euler maps, and using contraction of $\Pi_{p_t}$,
therefore yields, uniformly in $t$,
\begin{equation}
 \frac{W_2((\mathrm{Id}+\epsilon b_t)_\#p_t,p_t)}{|\epsilon|}
 \le k(t)+2d_R+C_R|\epsilon|
 \label{eq:uniform-tangent-upper}
\end{equation}
for $|\epsilon|\le1$, where $C_R$ is finite.  First let
$\epsilon\to0$ and then $R\to\infty$ for the uniform upper bound.

For the lower bound, write
$d_{t,\epsilon}=W_2((\mathrm{Id}+\epsilon b_t)_\#p_t,p_t)$.
The same-start coupling gives $d_{t,\epsilon}\le|\epsilon|M_b$.
Taylor expansion first in this coupling and then in an optimal
coupling gives, for every $\phi\in C_c^\infty$,
\begin{align}
 |\langle b_t,\nabla\phi\rangle_{p_t}|
 &\le \|\nabla\phi\|_{L^2(p_t)}
              \frac{d_{t,\epsilon}}{|\epsilon|}
       +|\epsilon|\|D^2\phi\|_\infty M_b^2.
 \label{eq:uniform-tangent-test-lower}
\end{align}
For each fixed test both its numerator and its gradient norm are
continuous in $t$: compose with $X_t(z)$ and use the stated continuity
and the integrable growth bound.  Fix $\delta>0$.  At each $t$ with
$k(t)\ge\delta$, the dual definition of projection supplies a test
with positive gradient norm and ratio at least $k(t)-\delta$.
Continuity gives a neighborhood in which that norm stays bounded
away from zero and its ratio is at least $k(s)-2\delta$.
The compact set $\{t:k(t)\ge\delta\}$ admits a finite subcover by
such neighborhoods.  Equation~\eqref{eq:uniform-tangent-test-lower}
for this finite collection consequently gives
$d_{t,\epsilon}/|\epsilon|\ge k(t)-2\delta-C_\delta|\epsilon|$
on this set.  On its complement nonnegativity gives a lower bound of
$k(t)-\delta$.  Let $\epsilon\to0$ and then $\delta\to0$.
\end{proof}

\subsection{The global trajectory expansion with only a second moment}

\begin{lemma}[Variational expansion of recursive Euler]
\label{lem:global-strong-expansion}
Under the grids and profiles of Theorem~\ref{thm:sharp-global-euler},
initialize Euler and the exact trajectory at the same $z$.  Then
\begin{equation}
 \left\|\sup_{0\le t\le T}
 \left|\frac{\widehat X_t^n-X_t}{\eta_n}
          +\frac12J_t C_t^\sigma\right|\right\|_{L^2(p_0)}
 \longrightarrow0.
 \label{eq:global-strong-expansion}
\end{equation}
\end{lemma}

\begin{proof}
Suppress $n$ temporarily, and write $e_k(z)=\widehat X_k(z)-X_{t_k}(z)$.
The exact local defect before its second-order expansion is
\[
 X_{t_{k+1}}-X_{t_k}-h_ku_{t_k}(X_{t_k})
 =\int_0^{h_k}(h_k-s)\acc_{t_k+s}(X_{t_k+s})\,\dd s.
\]
Equation~\eqref{eq:global-flow-growth} and discrete Gronwall therefore
give the pointwise estimate
\begin{equation}
 \max_k|e_k(z)|\le C H(1+|z|).
 \label{eq:global-pointwise-euler-bound}
\end{equation}
Taylor expansion along the exact trajectory also gives
\begin{align*}
 X_{t_{k+1}}={}&X_{t_k}+h_ku_{t_k}(X_{t_k})
      +\tfrac12h_k^2\acc_{t_k}(X_{t_k})+r_k,\\
 \|r_k\|_{L^2(p_0)}\le{}& C_jh_k^3/6.
\end{align*}
Put $D_k(z)=Du_{t_k}(X_{t_k}(z))$ and
\[
 n_k=u_{t_k}(X_{t_k}+e_k)-u_{t_k}(X_{t_k})-D_ke_k.
\]
The exact recurrence is
\begin{equation}
 e_{k+1}=(I+h_kD_k)e_k
       -\tfrac12h_k^2\acc_{t_k}(X_{t_k})+h_kn_k-r_k.
 \label{eq:global-error-recurrence}
\end{equation}
The last two terms are negligible after propagation and division by
$\eta_n$.  For the Taylor remainder, Minkowski's inequality gives
\[
 \frac1{\eta_n}\sum_k\|r_k\|_{L^2(p_0)}
 \le\frac{C_jTH_n^2}{6\eta_n}\longrightarrow0.
\]
For the nonlinear remainder, fix $z$.  All its exact and numerical
positions for small $H_n$ lie in one compact set depending on $z$.
On that set $Du$ is uniformly continuous in $(t,x)$.  If $\omega_z$
is a corresponding modulus in the spatial variable, then
\[
 \frac1{\eta_n}\sum_k h_k|n_k(z)|
 \le C\frac{H_n}{\eta_n}(1+|z|)
               \omega_z(CH_n(1+|z|))\longrightarrow0.
\]
The left side is bounded by $C'(1+|z|)$, since
$\|Du\|_{\rm op}\le L$ and $H_n/\eta_n$ is bounded.
Dominated convergence in $L^2(p_0)$ is therefore valid.  This is the
step that avoids a fourth-moment assumption or a global bound on
$D^2u$.

The fundamental matrix products of $I+h_kD_k(z)$ are bounded by
$e^{LT}$.  For each fixed $z$ they converge, uniformly over pairs of
grid nodes, to the fundamental matrix of
$\dot v=Du_t(X_t(z))v$, namely $J_t(z)J_s(z)^{-1}$.
Here is an explicit consistency justification: the coefficient is
continuous on $[0,T]$ and bounded by $L$; comparing one exact linear
step with $I+h_kD_k$ gives a defect bounded by
$C h_k\{\omega_{D,z}(H_n)+H_n\}$ in operator norm.
Summing and applying the uniform stability bound makes the total
error tend to zero.  Replacing the starting node $t_{k+1}$ of a
propagation product by $t_k$ adds $O(H_n)$ by the same flow bound.
Variation of constants in Eq.~\eqref{eq:global-error-recurrence}
now yields
\begin{equation}
 \left\|\max_k\left|
 \frac{e_k}{\eta_n}+\frac12J_{t_k}
       \sum_{j<k}\frac{h_{n,j}^2}{\eta_n}c_{t_{n,j}}
                  \right|\right\|_{L^2(p_0)}\longrightarrow0.
 \label{eq:global-discrete-variation}
\end{equation}
Indeed every pointwise matrix-replacement error tends to zero and is
bounded by $C(1+|z|)$ after summation, because
$\sum_j h_{n,j}^2/\eta_n\le TH_n/\eta_n$.  Thus these replacements
also pass to $L^2(p_0)$ by dominated convergence.

For each fixed $z$, continuity of $c_s(z)$ shows uniformly in the
upper grid node that
\[
 \sum_{j<k}\frac{h_{n,j}^2}{\eta_n}c_{t_{n,j}}(z)
 -\int_0^{t_{n,k}}\sigma_n(s)c_s(z)\,\dd s\longrightarrow0.
\]
The difference is bounded by
$(TH_n/\eta_n)\sup_{|s-r|\le H_n}|c_s(z)-c_r(z)|$.
Weak-star convergence of $\sigma_n$ gives, for every fixed $z,t$,
\[
 \int_0^t\sigma_n(s)c_s(z)\,\dd s
 \longrightarrow\int_0^t\sigma(s)c_s(z)\,\dd s.
\]
The primitives are equi-Lipschitz in $t$ for this fixed $z$, so the
convergence is uniform over $t\in[0,T]$.  Both differences are
dominated by $C(1+|z|)$, allowing one more application of dominated
convergence in $L^2(p_0)$.  This proves the node version of
Eq.~\eqref{eq:global-strong-expansion}.

Finally, Euler's continuous interpolation is the linear interpolation
of its node values.  The difference between the chord of the exact
trajectory and that trajectory is bounded by
$CH_n^2(1+|z|)$, using its second derivative $\acc_t(X_t(z))$.
After division by $\eta_n$ this vanishes in $L^2$.  The continuous
field $t\mapsto J_t(z)C_t^\sigma(z)$ has a uniformly vanishing
interpolation modulus for each fixed $z$, dominated by $C(1+|z|)$.
These two facts extend the node result uniformly to all times.
\end{proof}

\begin{proof}[Proof of Theorem~\ref{thm:sharp-global-euler}]
Let $b_t=-B_t^\sigma/2$.  Equations~\eqref{eq:global-reference-defect}
and \eqref{eq:global-strong-expansion}, using the common starting point
as a coupling, imply
\[
 \sup_tW_2\bigl(\widehat p_t^n,
              (\mathrm{Id}+\eta_nb_t)_\#p_t\bigr)=o(\eta_n).
\]
The regularity and growth properties of $b_t$ were proved above, and
Lemma~\ref{lem:global-coefficient-continuity} gives continuity of its
projected norm.  Apply Lemma~\ref{lem:uniform-global-tangent} and the
triangle inequality to obtain Eq.~\eqref{eq:sharp-global-uniform}.
Since $H_n\to0$, the maximum of the continuous coefficient over grid
nodes converges to its maximum over $[0,T]$.  This proves
Eq.~\eqref{eq:sharp-global-maximum}.
\end{proof}

\subsection{Quantile schedules and the characterization theorem}

For a positive $C^1$ function $\rho$ with integral one, the quantile
nodes obey, by the mean value theorem,
$Nh_k=1/\rho(\xi_k)$ for some $\xi_k\in[t_k,t_{k+1}]$.
The positive minimum of $\rho$ gives $H_N=O(N^{-1})$, and uniform
continuity gives $\sigma_N=Nh_k\to1/\rho$ uniformly on the intervals.
Taking $t=T$ in Theorem~\ref{thm:sharp-global-euler} proves
Eq.~\eqref{eq:sharp-schedule-coefficient}.  The projection can be moved
inside the integral because it is a bounded linear operator on the
fixed Hilbert space $L^2(p_T)$.  Also
$g_s^T=\Pi_{p_T}A_{s,T}$ is continuous in that space: after composition
with $X_T$, $A_{s,T}$ is $J_Tc_s$, a continuous $L^2(p_0)$ curve.

\begin{proof}[Proof of Theorem~\ref{thm:global-ap-characterization}]
Item 1 implies item 2 by Theorem~\ref{thm:distribution_order}; item 2
implies item 3 immediately.  For item 3, take $\eta_n=h_n=T/N_n$,
so $\sigma_n=1$.  Equation~\eqref{eq:sharp-global-maximum} forces
$B_t^1\in\mathcal G_{p_t}^\perp$ at every $t$.
Lemma~\ref{lem:global-solenoidal-transport} says equivalently
\[
 C_t^1=\int_0^tc_s\,\dd s\in S=\mathcal G_{p_0}^\perp
 \quad\text{for every }t.
\]
The curve is continuously differentiable in $L^2(p_0)$ with derivative
$c_t$.  Difference quotients lie in the closed subspace $S$, so
$c_t\in S$ for all $t$, including the endpoints by one-sided limits.
Transporting back gives
$\acc_t\in\mathcal G_{p_t}^\perp$, which is item 1.

Item 1 also implies item 4, because each fixed positive allocation
density has $H_N=O(N^{-1})$ and item 2 applies.  Conversely suppose
item 4 holds, even only for densities in a prescribed $C^1$
neighborhood of $1/T$.  Equation~\eqref{eq:sharp-schedule-coefficient}
gives $\int g_s^T/\rho(s)\,\dd s=0$ for each such density.
For an arbitrary smooth real function $f$, set
\[
 \sigma_\lambda(s)=1+\lambda f(s),\qquad
 Z_\lambda=\int_0^T\frac{\dd s}{\sigma_\lambda(s)},\qquad
 \rho_\lambda=\frac{\sigma_\lambda^{-1}}{Z_\lambda}.
\]
For sufficiently small $|\lambda|$, these are positive smooth
probability densities and $\rho_\lambda\to1/T$ in $C^1$.
Since $1/\rho_\lambda=Z_\lambda(1+\lambda f)$, subtracting the
$\lambda=0$ equality gives
$\int_0^Tf(s)g_s^T\,\dd s=0$.  Testing against every vector in
$L^2(p_T)$ and then using continuity in $s$ gives $g_s^T=0$ for
every $s$.  Thus $A_{s,T}$ is solenoidal under $p_T$; the inverse
of the transport in Lemma~\ref{lem:global-solenoidal-transport}
makes $\acc_s$ solenoidal under $p_s$.  This is item 1.
\end{proof}

For clarity, the positive first-order constant in the non-preserving
case on uniform grids is
\begin{equation}
 c_u=\frac12\max_{0\le t\le T}
 \left\|\Pi_{p_t}\int_0^t A_{s,t}\,\dd s\right\|_{L^2(p_t)}.
 \label{eq:global-positive-constant}
\end{equation}
If it vanished, the preceding proof would establish the acceleration
condition.  It is finite by the flow bounds and strictly positive
otherwise.  Equation~\eqref{eq:sharp-global-maximum} therefore gives
$\max_kW_2=c_uh+o(h)$, which proves the claimed $\Theta(h)$ conclusion.

If the uniform terminal coefficient vanishes but the acceleration
condition fails, then $\int_0^Tg_s^T\,\dd s=0$ but the continuous
Hilbert-valued curve $g_s^T$ is not identically zero.  Choose a smooth
$f$ such that $V=\int_0^Tf(s)g_s^T\,\dd s\ne0$.  The densities
$\rho_\lambda$ just constructed have exact terminal coefficient
$Z_\lambda|\lambda|\|V\|/2>0$ for every sufficiently small nonzero
$\lambda$.  This proves the stated sensitivity to arbitrarily small
schedule changes.
The same conclusion holds near any positive smooth density $\rho_0$
whose terminal coefficient vanishes: replace $1+\lambda f$ by
$1/\rho_0+\lambda f$ and normalize its reciprocal.  At $\lambda=0$
the normalizing integral is $\int\rho_0=1$, and the first-order
coefficient is again a positive scalar times $|\lambda|\|V\|/2$.

\begin{proof}[Proof of Corollary~\ref{cor:odd-global-order}]
Theorem~\ref{thm:nonlinear-odd-rigidity} excludes the acceleration
condition for every generator in the stated class when $m$ or $g$ is
nonaffine.  Equation~\eqref{eq:global-positive-constant} and
Theorem~\ref{thm:sharp-global-euler} then give the strictly positive
first-order constant for each fixed generator.  When both functions
are affine, the same rigidity theorem identifies the unique
acceleration-preserving generator and proves its acceleration is
identically zero.  Its exact trajectories are straight with constant
velocity, so Euler is exact on every grid.  Every other fixed
generator fails the acceleration condition and is covered by the
same positive-constant conclusion.
\end{proof}

\subsection{Finite schedule tests and the Gaussian counterexample}

Fix $T=1$, $a>0$, and a smooth real function $f$ with
$\|f\|_\infty<a^2$.  Put $b=\sqrt{a^2-f}$ and
$M_t=aI+b(t)J$.  The flow of $u_t(x)=M_tx$ is
\[
 F_{s,t}=e^{a(t-s)}
       \exp\left(J\int_s^tb(r)\,\dd r\right).
\]
It therefore generates $p_t=\mathcal N(0,e^{2at}I_2)$ from the
standard Gaussian.  Direct differentiation gives
\begin{equation}
 \acc_t(x)=\{f(t)I+[b'(t)+2ab(t)]J\}x,
 \qquad \Pi_{p_t}\acc_t=f(t)x.
 \label{eq:cosine-acceleration-projection}
\end{equation}
Indeed $Jx$ is solenoidal for every circular Gaussian, whereas $x$
is in its gradient closure by smooth cutoff approximation of
$|x|^2/2$.  All hypotheses of Theorem~\ref{thm:projected-residual}
hold: the fields are smooth and linear in $x$ with bounded matrix
coefficients on the compact time interval, and Gaussian moments are
finite.  The acceleration condition holds exactly when $f=0$.

Every Euler update is a scalar expansion followed by a rotation.
Writing its isotropic standard deviation after $k$ steps as $r_k$,
we have the exact recurrences and distance
\begin{align}
 r_0&=1,&
 r_{k+1}&=r_k\sqrt{(1+ah_k)^2+b(t_k)^2h_k^2},\nonumber\\
 W_2(\widehat p_k,p_{t_k})&=\sqrt2\,|r_k-e^{at_k}|.
 \label{eq:cosine-exact-euler-distance}
\end{align}
For sufficiently small $H$, uniformly in $t$ and $0\le h\le H$,
\begin{align}
 &\frac12\log\{1+2ah+(2a^2-f(t))h^2\}-ah\nonumber\\
 &\quad=-\frac12f(t)h^2
       +\left(-\frac23a^3+af(t)\right)h^3\nonumber\\
 &\qquad\quad+\left(a^4-a^2f(t)-\frac14f(t)^2\right)h^4
       +O(h^5).
 \label{eq:cosine-log-expansion}
\end{align}
The remainder is uniform because $f$ is bounded and the logarithm's
argument stays in a fixed neighborhood of one.  In particular, for
any refining grid,
\begin{equation}
 \log(r_N/e^a)=-\frac12\sum_kf(t_k)h_k^2
                     +O\left(\sum_kh_k^3\right).
 \label{eq:finite-schedule-log-error}
\end{equation}

For a fixed positive $C^1$ density $\rho$, its quantile grid satisfies
\begin{equation}
 \sum_kf(t_k)h_k^2
 =\frac1N\int_0^1\frac{f(s)}{\rho(s)}\,\dd s+O(N^{-2}).
 \label{eq:weighted-quantile-riemann}
\end{equation}
To see the order of the remainder, write the sum as
$N^{-1}\sum_k f(t_k)(Nh_k)h_k$.  The mean value formula and the
$C^1$ bound on $1/\rho$ give
$Nh_k=1/\rho(t_k)+O(H_N)$ uniformly.  The remaining left Riemann
sum for $f/\rho$ has error $O(H_N)$, and $H_N=O(N^{-1})$.
Also $\sum h_k^3=O(N^{-2})$.  Thus vanishing of the integral in
Eq.~\eqref{eq:weighted-quantile-riemann} gives an actual
$O(N^{-2})$ endpoint error through
Eq.~\eqref{eq:cosine-exact-euler-distance}, rather than merely the
$o(N^{-1})$ conclusion supplied by a zero first-order coefficient.

\begin{proof}[Proof of Proposition~\ref{prop:finite-schedule-blindness}]
Choose $m+1$ nonzero smooth bump functions with disjoint compact
supports in $(0,1)$.  The $m$ constraints
$\int_0^1 f/\rho_i=0$ form a homogeneous linear system in their
$m+1$ coefficients, so there is a nonzero solution.  Disjoint
supports make the resulting $f$ nonzero.  Multiplying it by a
sufficiently small nonzero scalar gives $\|f\|_\infty<a^2$ and
arbitrarily small size in any fixed finite number of derivative
norms.  The same is true of $\sqrt{a^2-f}-a$ by smoothness of the
square root near $a^2$.  Equations~\eqref{eq:finite-schedule-log-error}
and \eqref{eq:weighted-quantile-riemann} prove all the endpoint
$O(N^{-2})$ bounds.  Equation~\eqref{eq:cosine-acceleration-projection}
shows failure of the acceleration condition, and
Theorem~\ref{thm:global-ap-characterization} gives
$\Theta(N^{-1})$ maximum error on uniform grids.

For finitely many observation times $\tau_j$, impose the additional
linear constraints $\int_0^{\tau_j}f/\rho_i=0$ and choose more
disjoint bumps than constraints.  If the time is not a grid node,
insert it into the quantile grid or use the Euler continuous
interpolation.  Splitting the one interval containing a fixed time
changes the sum in Eq.~\eqref{eq:finite-schedule-log-error} by
$O(N^{-2})$.  The same endpoint argument at $\tau_j$ applies.
This is a finite collection of observation times, not a dense family
whose size grows with $N$.
\end{proof}

\subsection{The explicit cosine example and its exact coefficients}

Take $f(t)=\epsilon\cos(2\pi t)$, with $0<\epsilon<a^2$.
For uniform steps $h=1/N$ and $N\ge3$, discrete Fourier sums give
\[
 \sum_{k=0}^{N-1}f(k/N)=0,\qquad
 \sum_{k=0}^{N-1}f(k/N)^2=N\epsilon^2/2.
\]
Summing Eq.~\eqref{eq:cosine-log-expansion} gives the stronger expansion
\begin{equation}
 \log(r_N/e^a)
 =-\frac{2a^3}{3N^2}
    +\frac{a^4-\epsilon^2/8}{N^3}+O(N^{-4}).
 \label{eq:cosine-terminal-log-expansion}
\end{equation}
Since $W_2=\sqrt2 e^a|\exp(\log(r_N/e^a))-1|$, this proves
Eq.~\eqref{eq:cosine-terminal-second}.  Setting $\epsilon=0$
gives the preserving field $a(I+J)x$ and exactly the same leading
second-order coefficient.

Because all involved matrices commute with the rotation flow,
\[
 A_{s,t}(x)=\{f(s)I+[b'(s)+2ab(s)]J\}x.
\]
Projection under $p_t$ removes its $Jx$ component.  Thus the entire
first-order error profile for uniform grids, including continuous
interpolation, is
\begin{equation}
 \lim_{N\to\infty}N W_2(\widehat p_t^N,p_t)
 =\frac{\sqrt2\,e^{at}\epsilon}{4\pi}|\sin(2\pi t)|,
 \label{eq:cosine-uniform-first-profile}
\end{equation}
uniformly in $t$.  At $t=1/4$, taking $N$ divisible by four
identifies the interpolation with the numerical node and proves the
displayed intermediate-time coefficient in the main text.  The
maximum coefficient is positive and, if an explicit value is desired,
its maximizer and value are
\begin{align*}
 t_*&=1-\frac{1}{2\pi}\arctan(2\pi/a),\\
 C_{\max}&=\frac{\sqrt2\,\epsilon}{4\pi}
    e^{at_*}\sin\!\bigl(\arctan(2\pi/a)\bigr).
\end{align*}
There is one local maximum in each half interval; the second is
larger by the factor $e^{a/2}>1$.

Finally let $\rho_\delta(t)=1+\delta\cos(2\pi t)$ with
$0<|\delta|<1$.  Its quantile nodes are the unique solutions of
\[
 t_k+\frac{\delta}{2\pi}\sin(2\pi t_k)=k/N.
\]
The classical integral
$\int_0^1(1+\delta\cos(2\pi t))^{-1}\,\dd t
=(1-\delta^2)^{-1/2}$ follows, for example, by the substitution
$z=\tan(\pi t)$ on the two half intervals.  Since
$\cos\theta/(1+\delta\cos\theta)
=\delta^{-1}\{1-(1+\delta\cos\theta)^{-1}\}$,
\[
 \int_0^1\frac{f(t)}{\rho_\delta(t)}\,\dd t
 =\frac{\epsilon}{\delta}
          \left(1-\frac1{\sqrt{1-\delta^2}}\right).
\]
Equation~\eqref{eq:sharp-schedule-coefficient}, or directly
Eqs.~\eqref{eq:finite-schedule-log-error}--\eqref{eq:weighted-quantile-riemann},
now gives Eq.~\eqref{eq:cosine-perturbed-first}.  Its coefficient is
strictly positive for every nonzero $\delta$, however small.
